\section{The Thalean Quadratic}

\subsection{The Petersen-rooted sector system}

Let
\[
\Gfifteen\cong L(\Petersen)
\]
and let \(A\) be its adjacency matrix.  The graph has fifteen
vertices and thirty edges.

For each \(v\in V(\Gfifteen)\), the transport construction supplies a
sector
\[
S(v)\subseteq E(\Gfifteen).
\]
Every sector contains fourteen edges, and every core edge belongs to
exactly seven sectors.

Define
\[
M=(M_{v,e})\in\{0,1\}^{15\times30}
\]
by
\[
M_{v,e}
=
\begin{cases}
1,&e\in S(v),\\
0,&e\notin S(v).
\end{cases}
\]

The Thalean quadratic is
\[
Q:=MM^{\mathsf T}.
\]
Its entries are
\[
Q_{u,v}
=
|S(u)\cap S(v)|.
\]

\subsection{Distance and adjacency forms}

The overlap depends only on distance in \(\Gfifteen\):
\[
Q
=
14D_0+9D_1+5D_2+4D_3.
\]
Because \(\Gfifteen\) is distance-regular of diameter three, this
reduces exactly to
\[
Q=A^3+2A^2+2I.
\]

The identity has two logically separate interpretations.

First, it is a Gram identity:
\[
Q=MM^{\mathsf T}.
\]
Second, it belongs to the adjacency algebra:
\[
Q=A^3+2A^2+2I.
\]

Neither description includes the signed lift or its holonomy class.

\subsection{Edge-mediated accumulation}

For a sector coefficient vector
\[
x\in\mathbb R^{15},
\]
the vector
\[
M^{\mathsf T}x\in\mathbb R^{30}
\]
assigns to each core edge the sum of the coefficients of all sectors
containing that edge.

Applying \(M\) again returns the accumulated edge values to the
sector register:
\[
x
\longmapsto
M^{\mathsf T}x
\longmapsto
MM^{\mathsf T}x
=
Qx.
\]

Thus \(Q\) is an exact finite local-to-global-to-local incidence sum.

\subsection{The common incidence mode}

Let \(\one_{15}\) and \(\one_{30}\) denote the all-ones vectors in
their respective dimensions.  Column weight seven gives
\[
M^{\mathsf T}\one_{15}
=
7\one_{30}.
\]
Row weight fourteen gives
\[
M\one_{30}
=
14\one_{15}.
\]
Consequently,
\[
Q\one_{15}
=
M(7\one_{30})
=
98\one_{15}.
\]

The number \(98=7\cdot14\) is therefore obtained directly from the
balanced incidence structure.

\subsection{Spectrum}

The adjacency spectrum is
\[
\spec(A)
=
\left\{
4^{(1)},\,
2^{(5)},\,
(-1)^{(4)},\,
(-2)^{(5)}
\right\}.
\]
For
\[
p(t)=t^3+2t^2+2,
\]
we have
\[
p(4)=98,\qquad
p(2)=18,\qquad
p(-1)=3,\qquad
p(-2)=2.
\]
Hence
\[
\spec(Q)
=
\left\{
98^{(1)},\,
18^{(5)},\,
3^{(4)},\,
2^{(5)}
\right\}.
\]

In particular, \(Q\) is symmetric positive definite.

\subsection{Canonical normalization}

Define
\[
\widehat{Q}
:=
\frac{1}{98}Q.
\]
Then
\[
\widehat{Q}\one_{15}
=
\one_{15}.
\]

Since \(Q\) is symmetric and every row sum is \(98\),
\(\widehat{Q}\) is symmetric and doubly stochastic.  No probabilistic
interpretation is required.  Here it is used only as the unique scalar
normalization that fixes the common incidence mode.

Its spectrum is
\[
\spec(\widehat{Q})
=
\left\{
1^{(1)},\,
\left(\frac{9}{49}\right)^{(5)},\,
\left(\frac{3}{98}\right)^{(4)},\,
\left(\frac{1}{49}\right)^{(5)}
\right\}.
\]

\subsection{Operator boundary}

The map \(\widehat{Q}\) is a canonically normalized incidence
operator.  The present theorem does not establish that one
application of \(\widehat{Q}\) is:

\begin{itemize}
\item one native transport step;
\item one Thalean-time increment;
\item one physical averaging event;
\item or one stage of physical compression.
\end{itemize}

Iteration is studied algebraically.  Its promotion to native loading
requires a separate crosswalk.
